\documentclass[12pt]{article}
\usepackage{amsmath,amsfonts,amssymb,amsthm}
\usepackage{booktabs}
\usepackage{enumitem}
\usepackage{hyperref}
\usepackage{lineno}
\usepackage{geometry}
\geometry{margin=2.5cm}

%% For final submission, replace this preamble with:
%% \documentclass[review,12pt]{elsarticle}
%% and use \begin{frontmatter}...\end{frontmatter}

\modulolinenumbers[5]

\newtheorem{theorem}{Theorem}
\newtheorem{lemma}[theorem]{Lemma}
\newtheorem{proposition}[theorem]{Proposition}
\newtheorem{corollary}[theorem]{Corollary}
\newtheorem{definition}[theorem]{Definition}
\newtheorem{assumption}{Assumption}
\newtheorem{remark}{Remark}
\newtheorem{counterexample}{Counterexample}

\newcommand{\R}{\mathbb{R}}
\newcommand{\E}{\mathbb{E}}
\newcommand{\MLP}{\mathrm{MLP}}
\newcommand{\norm}[1]{\left\|#1\right\|}
\newcommand{\abs}[1]{\left|#1\right|}
\newcommand{\Rad}{\mathfrak{R}}

\begin{document}

\title{Structural Inductive Bias Reduces Sample Complexity\\
in Implicit Neural Models for Dynamical Systems}

\author{Rodolfo O.\ Rodrigo\\[4pt]
  \small Instituto de Autom\'atica (INAUT),
  CONICET--Universidad Nacional de San Juan, Argentina\\
  \small \texttt{rrodrigo@inaut.unsj.edu.ar}}
\date{}
\maketitle

\begin{abstract}
We prove that implicit neural models with structural inductive bias
achieve tighter generalization bounds than explicit models, and
therefore require fewer training samples for the same approximation
quality.
The key mechanism is that an implicit formulation
$F(x, y, \MLP(x;\theta)) = 0$, where $F$ encodes known structural
knowledge about the target function, \emph{changes} the learning
target: the network approximates a simpler residual $h^*$ instead
of the original function $f^*$.
Using Rademacher complexity theory, we derive explicit generalization
bounds for both explicit and implicit formulations, and establish a
quantitative gain condition under which the implicit model achieves
lower sample complexity.
We define the \emph{structural content} of $F$ as the ratio of path
norm reduction, and show that the implicit model outperforms the
explicit one if and only if the structural content exceeds an
amplification factor determined by the Lipschitz constant of $F$.
We provide sufficient conditions for complexity reduction (structural
factorization, nonlinearity absorption), counterexamples showing
when implicit models fail to help, and two instantiations of the
framework:
physics-informed variable transformations, where $F$ absorbs known
solution structure such as oscillatory or exponential behavior,
reducing the learning target to a slowly varying modulation; and
the Homotopy Analysis Method (HAM), where $F$ encodes the governing
differential equation and decomposes the solution into a series of
linear corrections of decreasing complexity.
We distinguish this target-transformation mechanism from the
regularization approach of Physics-Informed Neural Networks (PINNs),
which constrains the optimization landscape but does not change the
learning target.
The connection to a deeper topological framework based on
autosimilar function spaces is identified as a direction for future
work.
\end{abstract}

\noindent\textbf{Keywords:} Sample complexity; Rademacher complexity;
structural inductive bias; implicit neural models;
Homotopy Analysis Method; physics-informed architectures

\linenumbers

%% =========================================================
\section{Introduction}
\label{sec:intro}
%% =========================================================

Neural networks have become the dominant tool for approximating
unknown functions from data, with universal approximation theorems
guaranteeing that sufficiently large networks can approximate any
continuous function on a compact domain
\cite{Cybenko1989, Hornik1991, Funahashi1989}.
However, universality is a statement about \emph{capacity}, not about
\emph{efficiency}: knowing that a network \emph{can} approximate
$f^*$ says nothing about how many parameters or training samples are
needed to do so reliably.
The gap between capacity and efficiency is governed by the
\emph{complexity} of the hypothesis class relative to the target
function.

In the context of dynamical systems, the target function $f^*$ is
not arbitrary: it is constrained by physical laws, symmetries, and
structural regularities.
A neural network that ignores these constraints must discover all
structure from data alone.
A network that exploits them --- by embedding known structure into
the model --- should, in principle, require less data for the same
approximation quality.
This intuition is widely held but rarely formalized.

Several strategies have been proposed for injecting physical
knowledge into neural models.
Physics-Informed Neural Networks (PINNs) \cite{Raissi2019} add a
PDE residual term to the loss function, penalizing solutions
incompatible with the governing equation.
Neural ODEs \cite{Chen2018} parametrize the vector field
$\dot{x} = f_\theta(x)$ and integrate it with a standard solver.
Deep Equilibrium Models (DEQs) \cite{Bai2019} define the output as
a fixed point of an implicit layer.
These approaches constrain the optimization landscape but share a
common feature: the network still approximates the full target
function $f^*$ directly.
The physics acts as a regularizer, not as a simplifier of the
learning target.

A fundamentally different strategy is to \emph{transform} the
learning target by encoding structural knowledge in the model
architecture.
If $f^*(x) = g(x) \cdot e^{-\lambda x}$ with $\lambda$ known, one
can define $\hat{f}(x) = \MLP(x;\theta) \cdot e^{-\lambda x}$, so
the network learns only the slowly varying modulation $g$ instead of
the full product $f^*$.
The known exponential structure is \emph{absorbed} by the model; the
network's learning target is genuinely simpler.

The present paper formalizes this target-transformation mechanism.
We prove that when a known function $F$ encodes structural knowledge
about $f^*$ and the network participates in an implicit equation
$F(x, y, \MLP(x;\theta)) = 0$, the learning target changes from
$f^*$ to a simpler residual $h^*$, and the sample complexity
decreases quantifiably.

\begin{remark}[Distinction from PINNs]
\label{rem:pinn_distinction}
It is important to distinguish the mechanism formalized here from the
PINN approach \cite{Raissi2019}.
In a PINN, the network $u_\theta(x)$ approximates $f^*$ directly;
the PDE residual $\mathcal{N}[u_\theta]$ is added to the loss as a
regularizer.
The learning target remains $f^* $--- the network must represent the
full solution.
In our framework, the implicit function $F$ \emph{changes} the
learning target from $f^*$ to a different, simpler function
$h^*$ defined by $F(x, f^*(x), h^*(x)) = 0$.
The network approximates $h^*$, which has lower complexity.

PINNs reduce effective complexity by constraining the optimization
landscape (a property of the training algorithm).
The present framework reduces effective complexity by reducing the
\emph{intrinsic difficulty of the approximation target} (a property
of the function class), which translates into tighter generalization
bounds independently of the optimization method.
The two mechanisms are complementary, not competing: one could use a
PINN-style residual loss \emph{in addition to} a target-transforming
implicit model.
\end{remark}

\subsection{Contributions}

\begin{enumerate}
  \item We derive explicit Rademacher complexity bounds for both
    explicit and implicit formulations
    (Theorem~\ref{teo:principal}), with a rigorous treatment of the
    Lipschitz composition and centering correction required by the
    implicit model.
  \item We establish a quantitative \emph{gain condition}
    (Corollary~\ref{cor:datos}) under which the implicit model
    requires fewer training samples, expressed in terms of the
    path norm reduction and the Lipschitz constant of~$F$.
  \item We define the \emph{structural content} of $F$
    (Definition~\ref{def:structural_content}) as the ratio of path
    norm bounds, and provide sufficient conditions for $F$ to have
    positive structural content (Propositions~\ref{prop:factorizacion}
    and~\ref{prop:absorcion}).
  \item We show that physics-informed variable transformations and
    HAM-based decompositions are instances of this framework
    (Section~\ref{sec:instances}), explaining why architectures
    such as SIREN \cite{Sitzmann2020} and Laguerre networks
    outperform generic MLPs for specific dynamical classes.
  \item We provide counterexamples (Section~\ref{sec:counter})
    showing that not every implicit formulation reduces complexity,
    and characterize the conditions under which the implicit model
    is actually worse.
\end{enumerate}

The connection between the structural inductive bias formalized here
and a deeper topological framework based on autosimilar function
spaces and iterated function systems is identified as a direction for
future work (Section~\ref{sec:discussion}).

%% =========================================================
\section{Problem Formulation}
\label{sec:formulation}
%% =========================================================

\begin{definition}[Learning scenario]
\label{def:scenario}
Let $\mathcal{X} \subset \R^n$ be compact with $\norm{x} \le X_{\max}$
for all $x \in \mathcal{X}$.
Let $f^*: \mathcal{X} \to \R$ be the unknown target function.
A training set $S = \{(x_i, y_i)\}_{i=1}^N$ is available with
$y_i = f^*(x_i) + \varepsilon_i$, where $\varepsilon_i$ is
sub-Gaussian noise with parameter $\sigma^2$.
\end{definition}

\begin{definition}[Hypothesis class]
\label{def:hypothesis}
For an MLP of depth $L$ and maximum width $d$, the hypothesis class
with bounded spectral path norm is:
\begin{equation}
  \mathcal{H}_B = \left\{ \MLP(\cdot\,;\theta) :
    \prod_{l=1}^{L} \norm{W_l}_\sigma \le B \right\}
\end{equation}
where $\norm{W_l}_\sigma$ is the spectral norm of the weight matrix
at layer $l$ and $B > 0$ is the path norm bound.
\end{definition}

We consider two formulations for learning $f^*$:

\textbf{Explicit formulation.}
The network approximates $f^*$ directly:
\begin{equation}
  \hat{\theta}_{\mathrm{exp}} = \arg\min_\theta \;
  \hat{\mathcal{L}}_{\mathrm{exp}}(\theta)
  = \arg\min_\theta \;
  \frac{1}{N}\sum_{i=1}^{N}(y_i - \MLP(x_i;\theta))^2.
  \label{eq:loss_exp}
\end{equation}

\textbf{Implicit formulation.}
Let $F: \mathcal{X} \times \R \times \R \to \R$ be a known function
of class $C^1$.
The network participates as a component of an implicit equation:
\begin{equation}
  \hat{\theta}_{\mathrm{imp}} = \arg\min_\theta \;
  \hat{\mathcal{L}}_{\mathrm{imp}}(\theta)
  = \arg\min_\theta \;
  \frac{1}{N}\sum_{i=1}^{N} F(x_i, y_i, \MLP(x_i;\theta))^2.
  \label{eq:loss_imp}
\end{equation}

The function $F$ encodes structural knowledge about the relationship
between $x$ and $y$.
When $F(x,y,z) = y - z$, the implicit formulation reduces to the
explicit one.
When $F$ is nontrivial, it absorbs part of the complexity of $f^*$,
so that the network learns a different, simpler function
$h^* \ne f^*$.

%% =========================================================
\section{Formal Hypotheses}
\label{sec:hypotheses}
%% =========================================================

\begin{assumption}[Regularity of $F$]
\label{hip:F}
The function $F: \mathcal{X} \times \R \times \R \to \R$ satisfies:
\begin{enumerate}[label=(\alph*)]
  \item $F$ is $C^1$ in all variables.
  \item \textbf{Local Lipschitz in the training region.}
    Let $\mathcal{R} = \mathcal{X} \times \mathcal{Y} \times
    \mathcal{Z}$ be the training region, where
    $\mathcal{Y} \supset \{f^*(x) : x \in \mathcal{X}\}$ and
    $\mathcal{Z} \supset \{h^*(x) : x \in \mathcal{X}\}$ are compact
    neighborhoods.
    There exist constants $0 < \alpha \le \beta$ such that:
    \begin{equation}
      \alpha \;\le\; \abs{\frac{\partial F}{\partial z}(x,y,z)}
      \;\le\; \beta \qquad \forall\,(x,y,z) \in \mathcal{R}.
      \label{eq:sandwich}
    \end{equation}
    The constant $\beta$ is \emph{local}: it depends on the size of
    $\mathcal{R}$ and is not required to hold globally.
  \item $\frac{\partial F}{\partial y}(x,y,z) \neq 0$ on
    $\mathcal{R}$ (Implicit Function Theorem condition).
\end{enumerate}
\end{assumption}

\begin{assumption}[Existence of implicit target]
\label{hip:h}
Under Assumption~\ref{hip:F}, there exists a function
$h^*: \mathcal{X} \to \R$ of class $C^1$, defined implicitly by:
\begin{equation}
  F(x, f^*(x), h^*(x)) = 0 \qquad \forall\, x \in \mathcal{X}.
\end{equation}
The function $h^*$ is what the MLP must learn in the implicit
formulation.
\end{assumption}

\begin{assumption}[Complexity reduction]
\label{hip:reduccion}
There exist path norm bounds $B_{\mathrm{exp}}$ and $B_{\mathrm{imp}}$
such that:
\begin{enumerate}[label=(\alph*)]
  \item $f^*$ is $\varepsilon_0$-approximable by
    $\mathcal{H}_{B_{\mathrm{exp}}}$.
  \item $h^*$ is $\varepsilon_0$-approximable by
    $\mathcal{H}_{B_{\mathrm{imp}}}$.
  \item $B_{\mathrm{imp}} < B_{\mathrm{exp}}$.
\end{enumerate}
\end{assumption}

\begin{remark}[Nature of Assumption~\ref{hip:reduccion}]
\label{rem:nature}
Assumption~\ref{hip:reduccion} is the substantive hypothesis.
Assumptions~\ref{hip:F} and~\ref{hip:h} are technical regularity
conditions.
The claim ``the implicit model requires fewer data'' is exactly as
strong as Assumption~\ref{hip:reduccion} --- it depends entirely on
how much complexity $F$ absorbs.
Sections~\ref{sec:sufficient} and~\ref{sec:counter} provide
conditions under which this assumption holds or fails.
\end{remark}

%% =========================================================
\section{Preliminary Results}
\label{sec:prelim}
%% =========================================================

\begin{lemma}[Rademacher complexity of norm-bounded networks ---
Golowich, Rakhlin \& Shamir \cite{Golowich2018}]
\label{lema:rad}
For the class $\mathcal{H}_B$ of networks of depth $L$, maximum
width $d$, and spectral path norm bounded by $B$:
\begin{equation}
  \Rad_N(\mathcal{H}_B) \;\le\;
  \frac{B \cdot X_{\max} \cdot \sqrt{2L\log(2d)}}{\sqrt{N}}.
  \label{eq:rademacher_nn}
\end{equation}
\end{lemma}

\begin{remark}[On the simplification]
\label{rem:simplification}
The full bound of Golowich et al.\ involves per-layer spectral
norms and $(2,1)$-norms:
\begin{equation}
  \Rad_N(\mathcal{H}_B) \;\le\;
  \frac{\bigl(\prod_{l=1}^{L}\norm{W_l}_\sigma\bigr)
    \cdot X_{\max} \cdot \sqrt{\sum_{l=1}^{L}
    \bigl(\norm{W_l}_{2,1}/\norm{W_l}_\sigma\bigr)^{2/3}}}
  {\sqrt{N}}.
\end{equation}
Bound~\eqref{eq:rademacher_nn} follows from
$\norm{W_l}_{2,1}/\norm{W_l}_\sigma \le \sqrt{d}$.
This preserves the asymptotic order $O(1/\sqrt{N})$ but may be loose
for structurally sparse weight matrices.
\end{remark}

\begin{lemma}[Contraction lemma --- Ledoux \& Talagrand
\cite{Ledoux1991}]
\label{lema:contraccion}
Let $\phi: \R \to \R$ be $L_\phi$-Lipschitz with $\phi(0) = 0$,
and let $\mathcal{G}$ be a function class.
Then:
\begin{equation}
  \Rad_N(\phi \circ \mathcal{G})
  \;\le\; L_\phi \cdot \Rad_N(\mathcal{G}).
\end{equation}
\end{lemma}

\begin{lemma}[Uniform generalization bound --- Bartlett \& Mendelson
\cite{Bartlett2002}]
\label{lema:generalizacion}
Let $\mathcal{G}$ be a function class with values in $[0, M]$.
Then, with probability at least $1 - \delta$ over $S$:
\begin{equation}
  \sup_{g \in \mathcal{G}}
  \abs{\E[g] - \hat{\E}_N[g]}
  \;\le\; 2\Rad_N(\mathcal{G})
  + 3M\sqrt{\frac{\log(2/\delta)}{2N}}.
  \label{eq:generalizacion}
\end{equation}
\end{lemma}

%% =========================================================
\section{Main Theorem}
\label{sec:main}
%% =========================================================

\begin{theorem}[Generalization bound reduction by implicit model]
\label{teo:principal}
Under Assumptions~\ref{hip:F}--\ref{hip:reduccion}, let
$\hat{\theta}_{\mathrm{exp}} \in \mathcal{H}_{B_{\mathrm{exp}}}$ and
$\hat{\theta}_{\mathrm{imp}} \in \mathcal{H}_{B_{\mathrm{imp}}}$ be
the empirical risk minimizers of
$\hat{\mathcal{L}}_{\mathrm{exp}}$ and
$\hat{\mathcal{L}}_{\mathrm{imp}}$ respectively.
Then, with probability at least $1 - \delta$:

\medskip

\textbf{(I) Explicit bound:}
\begin{equation}
  \mathcal{L}_{\mathrm{exp}}(\hat{\theta}_{\mathrm{exp}})
  \;\le\; \hat{\mathcal{L}}_{\mathrm{exp}}(\hat{\theta}_{\mathrm{exp}})
  + \frac{4 M_{\mathrm{exp}} \cdot B_{\mathrm{exp}} \cdot X_{\max}
    \cdot \sqrt{2L\log(2d)}}{\sqrt{N}}
  + 3 M_{\mathrm{exp}}^2
    \sqrt{\frac{\log(2/\delta)}{2N}}
  \label{eq:cota_exp}
\end{equation}

\textbf{(II) Implicit bound:}
\begin{equation}
  \mathcal{L}_{\mathrm{imp}}(\hat{\theta}_{\mathrm{imp}})
  \;\le\; \hat{\mathcal{L}}_{\mathrm{imp}}(\hat{\theta}_{\mathrm{imp}})
  + \frac{4 \beta \cdot M_{\mathrm{imp}} \cdot B_{\mathrm{imp}}
    \cdot X_{\max} \cdot \sqrt{2L\log(2d)}}{\sqrt{N}}
  + 3 M_{\mathrm{imp}}^2
    \sqrt{\frac{\log(2/\delta)}{2N}}
  \label{eq:cota_imp}
\end{equation}

\noindent where
$M_{\mathrm{exp}} = \sup_{(x,\theta)} |y - \MLP(x;\theta)|$ and
$M_{\mathrm{imp}} = \sup_{(x,\theta)} |F(x,y,\MLP(x;\theta))|$ are
range bounds over $\mathcal{H}_{B_{\mathrm{exp}}}$ and
$\mathcal{H}_{B_{\mathrm{imp}}}$ respectively.
\end{theorem}

\begin{proof}
\textbf{Part~(I): Explicit bound.}

Define the loss class:
\begin{equation}
  \mathcal{G}_{\mathrm{exp}} = \{(x,y) \mapsto (y - g(x))^2
  : g \in \mathcal{H}_{B_{\mathrm{exp}}}\}.
\end{equation}

\emph{Step~1: Lipschitz composition.}
The function $\phi(t) = t^2$ is $2M_{\mathrm{exp}}$-Lipschitz on
$[-M_{\mathrm{exp}}, M_{\mathrm{exp}}]$.
The inner class
$\tilde{\mathcal{G}}_{\mathrm{exp}} = \{(x,y) \mapsto y - g(x)
: g \in \mathcal{H}_{B_{\mathrm{exp}}}\}$
satisfies
$\Rad_N(\tilde{\mathcal{G}}_{\mathrm{exp}})
= \Rad_N(\mathcal{H}_{B_{\mathrm{exp}}})$
(translation by $y$ does not affect Rademacher complexity, since $y$
is fixed with respect to the hypothesis class).

To apply the contraction lemma (Lemma~\ref{lema:contraccion}), we
note that $\phi(t) = t^2$ does not satisfy $\phi(0) = 0$.
However, the Rademacher complexity is invariant under constant shifts
of the function class: for any constant $c$,
$\Rad_N(\{g + c : g \in \mathcal{G}\}) = \Rad_N(\mathcal{G})$.
Since $(y - g(x))^2 = \psi(y - g(x))$ where $\psi(t) = t^2$ and
$\psi(0) = 0$, the contraction lemma applies to the centered class
$\tilde{\mathcal{G}}_{\mathrm{exp}}$ with Lipschitz constant
$2M_{\mathrm{exp}}$.
Therefore:
\begin{equation}
  \Rad_N(\mathcal{G}_{\mathrm{exp}})
  \;\le\; 2M_{\mathrm{exp}} \cdot
  \Rad_N(\mathcal{H}_{B_{\mathrm{exp}}}).
\end{equation}

\emph{Step~2: Rademacher substitution.}
By Lemma~\ref{lema:rad}:
\begin{equation}
  \Rad_N(\mathcal{G}_{\mathrm{exp}})
  \;\le\; \frac{2 M_{\mathrm{exp}} \cdot B_{\mathrm{exp}} \cdot
    X_{\max} \cdot \sqrt{2L\log(2d)}}{\sqrt{N}}.
\end{equation}

\emph{Step~3: Generalization bound.}
Applying Lemma~\ref{lema:generalizacion} yields~\eqref{eq:cota_exp}.
\hfill $\square_{\mathrm{I}}$

\medskip

\textbf{Part~(II): Implicit bound.}

Define:
\begin{equation}
  \mathcal{G}_{\mathrm{imp}} = \{(x,y) \mapsto F(x,y,g(x))^2
  : g \in \mathcal{H}_{B_{\mathrm{imp}}}\}.
\end{equation}

\emph{Step~1: Lipschitz composition with centering correction.}
The composition has two layers:
\begin{enumerate}[label=(\roman*)]
  \item $z \mapsto F(x,y,z)$ is $\beta$-Lipschitz in $z$ by
    Assumption~\ref{hip:F}(b).
  \item $t \mapsto t^2$ is $2M_{\mathrm{imp}}$-Lipschitz on
    $[-M_{\mathrm{imp}}, M_{\mathrm{imp}}]$.
\end{enumerate}

The contraction lemma (Lemma~\ref{lema:contraccion}) requires the
Lipschitz map to satisfy $\phi(0) = 0$.
The map $z \mapsto F(x,y,z)$ does not satisfy $F(x,y,0) = 0$ in
general.
We handle this by defining the centered map
$\tilde{F}_{x,y}(z) := F(x,y,z) - F(x,y,0)$, which satisfies
$\tilde{F}_{x,y}(0) = 0$ and is $\beta$-Lipschitz (same constant).
Since Rademacher complexity is invariant under addition of a function
that depends only on the data $(x,y)$ and not on the hypothesis $g$:
\begin{equation}
  \Rad_N\bigl(\{(x,y) \mapsto F(x,y,g(x))
    : g \in \mathcal{H}_{B_{\mathrm{imp}}}\}\bigr)
  = \Rad_N\bigl(\{(x,y) \mapsto \tilde{F}_{x,y}(g(x))
    : g \in \mathcal{H}_{B_{\mathrm{imp}}}\}\bigr).
\end{equation}

We now apply the contraction lemma twice.
First, to $\tilde{F}_{x,y}$ (which is $\beta$-Lipschitz with
$\tilde{F}_{x,y}(0) = 0$):
\begin{equation}
  \Rad_N\bigl(\{\tilde{F}_{x,y} \circ g
    : g \in \mathcal{H}_{B_{\mathrm{imp}}}\}\bigr)
  \;\le\; \beta \cdot \Rad_N(\mathcal{H}_{B_{\mathrm{imp}}}).
\end{equation}

Second, the squaring map $\psi(t) = t^2$ satisfies $\psi(0) = 0$
and is $2M_{\mathrm{imp}}$-Lipschitz on the relevant range.
Therefore:
\begin{equation}
  \Rad_N(\mathcal{G}_{\mathrm{imp}})
  \;\le\; 2M_{\mathrm{imp}} \cdot \beta \cdot
  \Rad_N(\mathcal{H}_{B_{\mathrm{imp}}}).
  \label{eq:rad_imp}
\end{equation}

\emph{Step~2: Rademacher substitution.}
By Lemma~\ref{lema:rad}:
\begin{equation}
  \Rad_N(\mathcal{G}_{\mathrm{imp}})
  \;\le\; \frac{2 \beta \cdot M_{\mathrm{imp}} \cdot
    B_{\mathrm{imp}} \cdot X_{\max} \cdot
    \sqrt{2L\log(2d)}}{\sqrt{N}}.
\end{equation}

\emph{Step~3: Generalization bound.}
Applying Lemma~\ref{lema:generalizacion}
yields~\eqref{eq:cota_imp}.
\hfill $\square_{\mathrm{II}}$
\end{proof}

\begin{remark}[Role of $\beta$]
\label{rem:beta}
The Lipschitz constant $\beta$ of $F$ in $z$ appears as an
amplification factor in the implicit bound~\eqref{eq:cota_imp}.
This is not an artifact of the proof: it reflects the genuine cost
of passing through the nonlinear map $F$.
If $F$ is ``rigid'' ($\beta$ large), the benefit of the reduced
path norm $B_{\mathrm{imp}} < B_{\mathrm{exp}}$ may be offset by
the amplification.
The gain condition is made precise in Corollary~\ref{cor:datos}.
\end{remark}

%% =========================================================
\section{Sample Complexity Comparison}
\label{sec:sample}
%% =========================================================

\begin{corollary}[Sample complexity reduction]
\label{cor:datos}
Fix a generalization gap $\varepsilon > 0$ and confidence
$1 - \delta$.
The number of training samples required to achieve the generalization
gap $\varepsilon$ is, up to logarithmic factors:

\textbf{Explicit:}
\begin{equation}
  N_{\mathrm{exp}} = O\!\left(
    \frac{M_{\mathrm{exp}}^2 \cdot B_{\mathrm{exp}}^2 \cdot
      X_{\max}^2 \cdot L\log(d)}{\varepsilon^2}
  \right).
  \label{eq:N_exp}
\end{equation}

\textbf{Implicit:}
\begin{equation}
  N_{\mathrm{imp}} = O\!\left(
    \frac{\beta^2 \cdot M_{\mathrm{imp}}^2 \cdot B_{\mathrm{imp}}^2
      \cdot X_{\max}^2 \cdot L\log(d)}{\varepsilon^2}
  \right).
  \label{eq:N_imp}
\end{equation}

The sample complexity ratio is:
\begin{equation}
  \boxed{
    \frac{N_{\mathrm{imp}}}{N_{\mathrm{exp}}}
    = \frac{\beta^2 \cdot M_{\mathrm{imp}}^2 \cdot
      B_{\mathrm{imp}}^2}
    {M_{\mathrm{exp}}^2 \cdot B_{\mathrm{exp}}^2}
  }
  \label{eq:ratio}
\end{equation}
\end{corollary}

\begin{proof}
Setting the dominant term in~\eqref{eq:cota_exp} equal to
$\varepsilon$ and solving for $N$ yields~\eqref{eq:N_exp}.
The procedure is identical for~\eqref{eq:N_imp}, with the additional
factor $\beta$ from the Lipschitz constant of $F$.
\end{proof}

\begin{remark}[Gain condition]
\label{rem:gain}
The implicit model achieves lower sample complexity when:
\begin{equation}
  \beta^2 \cdot M_{\mathrm{imp}}^2 \cdot B_{\mathrm{imp}}^2
  \;<\; M_{\mathrm{exp}}^2 \cdot B_{\mathrm{exp}}^2.
  \label{eq:condicion_reduccion}
\end{equation}
\end{remark}

\begin{definition}[Structural content of $F$]
\label{def:structural_content}
The structural content of $F$ is:
\begin{equation}
  \mathcal{S}(F) \;:=\; \frac{B_{\mathrm{exp}}}{B_{\mathrm{imp}}},
\end{equation}
with $\mathcal{S}(F) > 1$ indicating that $F$ provides useful
structural bias, and $\mathcal{S}(F) = 1$ indicating that $F$ is
structurally empty (cf.\ Counterexample~\ref{ce:trivial}).
The gain condition~\eqref{eq:condicion_reduccion} becomes:
\begin{equation}
  \mathcal{S}(F) \;>\;
  \frac{\beta \cdot M_{\mathrm{imp}}}{M_{\mathrm{exp}}}.
  \label{eq:condicion_S}
\end{equation}
\end{definition}

\begin{remark}[Structural content is problem-dependent]
\label{rem:problem_dependent}
The structural content $\mathcal{S}(F)$ depends on the pair
$(f^*, F)$, not on $F$ alone, since $B_{\mathrm{exp}}$ is determined
by the unknown $f^*$.
In practice, $\mathcal{S}(F)$ is not computable a priori.
However, the sufficient conditions of
Section~\ref{sec:sufficient} provide classes of problems for which
$\mathcal{S}(F) > 1$ can be guaranteed from structural properties
of $F$ alone.
\end{remark}

%% =========================================================
\section{Sufficient Conditions for Complexity Reduction}
\label{sec:sufficient}
%% =========================================================

Assumption~\ref{hip:reduccion} is the key piece.
Under what conditions on $F$ is
$B_{\mathrm{imp}} < B_{\mathrm{exp}}$ guaranteed?

\begin{proposition}[Reduction by structural factorization]
\label{prop:factorizacion}
Suppose $F$ has the form:
\begin{equation}
  F(x, y, z) = \psi(x, y) - z,
  \label{eq:F_factorizada}
\end{equation}
where $\psi: \mathcal{X} \times \R \to \R$ is a known function.
Then $h^*(x) = \psi(x, f^*(x))$, and the Lipschitz constant is
$\beta = 1$ (no amplification).

If $\psi$ simplifies the structure of $f^*$ --- for example, if
$f^*(x) = p(x) \cdot s(x)$ with $p$ a known polynomial and $s$ a
smooth function of low complexity, and
$\psi(x, y) = y / p(x)$ --- then $h^*(x) = s(x)$ and
$B_{\mathrm{imp}} \ll B_{\mathrm{exp}}$.
\end{proposition}

\begin{proof}
If $F(x,y,z) = \psi(x,y) - z$, then $F = 0$ implies
$z = \psi(x,y)$.
Evaluating at $y = f^*(x)$: $h^*(x) = \psi(x, f^*(x))$.
The Lipschitz constant is
$|\partial F / \partial z| = |-1| = 1$, so $\alpha = \beta = 1$.

For the path norm reduction: let $f^*(x) = p(x) \cdot s(x)$ where
$p$ is a known polynomial of degree $k$ and $s$ is smooth with slow
variation.
If $\psi(x,y) = y/p(x)$, then $h^*(x) = s(x)$.
By the Barron norm characterization of network approximation
complexity \cite{Barron1993}, the path norm required to approximate
a function is proportional to its spectral variation:
\begin{equation}
  B \;\propto\; C_f
  \;=\; \int_{\R^n} \norm{\omega} |\hat{f}(\omega)| \, d\omega.
\end{equation}

The product $f^* = p \cdot s$ corresponds to spectral convolution
$\widehat{p \cdot s} = \hat{p} * \hat{s}$, which redistributes the
spectral content of $s$ toward higher frequencies when $p$ is
nonconstant.
For polynomials of degree $k \ge 1$, the spectral measure $\hat{p}$
involves derivatives of the Dirac delta, ensuring redistribution
away from the origin.
Therefore $C_{f^*} = C_{p \cdot s} \ge C_s = C_{h^*}$ under generic
conditions on $p$, and $B_{\mathrm{imp}} \le B_{\mathrm{exp}}$.
\end{proof}

\begin{proposition}[Reduction by nonlinearity absorption]
\label{prop:absorcion}
Let $f^*(x) = \phi(g(x))$ where $\phi$ is a known invertible $C^1$
nonlinearity and $g$ is unknown.
Define:
\begin{equation}
  F(x, y, z) = \phi(z) - y.
\end{equation}
Then $h^*(x) = g(x) = \phi^{-1}(f^*(x))$, and the MLP only needs to
learn $g$ instead of $\phi \circ g$.

The Lipschitz constant is $\beta = \sup_{\mathcal{Z}} |\phi'(z)|$,
which may be large.
The gain condition~\eqref{eq:condicion_reduccion} is satisfied when
the reduction in $B$ dominates the amplification by $\beta$.
\end{proposition}

\begin{proof}
$F(x,y,z) = \phi(z) - y = 0$ implies $z = \phi^{-1}(y)$.
At $y = f^*(x) = \phi(g(x))$:
$h^*(x) = \phi^{-1}(\phi(g(x))) = g(x)$.

For $\phi(t) = e^t$: $f^* = e^g$ and $h^* = g = \ln f^*$.
The Barron norm of $e^g$ grows exponentially with that of $g$
(via the Taylor expansion of the exponential and iterated spectral
convolution), so $C_{f^*} = C_{e^g} \gg C_g = C_{h^*}$.

The Lipschitz constant is
$\beta = \sup |e^z| = e^{\sup |g|}$, which can be large.
The gain condition is met when the path norm reduction dominates:
$B_{\mathrm{exp}} / B_{\mathrm{imp}} > \beta \cdot
M_{\mathrm{imp}} / M_{\mathrm{exp}}$.
\end{proof}

%% =========================================================
\section{Counterexamples: When the Implicit Model Fails}
\label{sec:counter}
%% =========================================================

\begin{counterexample}[$F$ without structural content]
\label{ce:trivial}
Let $F(x,y,z) = y - z$.
Then $h^*(x) = f^*(x)$ and $B_{\mathrm{imp}} = B_{\mathrm{exp}}$.
The implicit equation absorbs no complexity; the network must learn
exactly the same function.
The Lipschitz constant is $\beta = 1$, so there is no amplification,
but also no gain: $\mathcal{S}(F) = 1$.
\end{counterexample}

\begin{counterexample}[$F$ that amplifies complexity]
\label{ce:amplification}
Let $F(x,y,z) = y - e^z$.
Then $h^*(x) = \ln(f^*(x))$, which may have singularities if $f^*$
vanishes, and the condition $\alpha > 0$ of Assumption~\ref{hip:F}
fails in regions where $e^z \to 0$.
Even when $f^* > 0$ everywhere, the Lipschitz constant
$\beta = \sup |e^z|$ can be large in the training region, potentially
making the implicit formulation worse than the explicit one.
\end{counterexample}

\begin{remark}[Diagnostic criterion]
\label{rem:diagnostic}
Not every implicit model reduces complexity.
The structural content $\mathcal{S}(F)$ is the diagnostic:
$\mathcal{S}(F) > 1$ is necessary, and the full gain
condition~\eqref{eq:condicion_S} is sufficient.
Counterexample~\ref{ce:trivial} shows $\mathcal{S} = 1$ (neutral),
and Counterexample~\ref{ce:amplification} shows how
the amplification factor $\beta$ can dominate even when $\mathcal{S}
> 1$.
\end{remark}

%% =========================================================
\section{Instantiations of the Framework}
\label{sec:instances}
%% =========================================================

We show that two strategies for injecting structural knowledge into
neural models are genuine instances of the general framework of
Theorem~\ref{teo:principal}: both change the learning target from
$f^*$ to a simpler function $h^*$.

\subsection{Physics-informed variable transformations}
\label{sec:variable_transform}

Consider a dynamical system whose solutions are known to exhibit a
specific qualitative structure: oscillatory behavior, exponential
decay, or polynomial modulation.
In each case, the practitioner possesses structural knowledge about
$f^*$ that can be absorbed into $F$, genuinely simplifying the
learning target.

\textbf{Case 1: Exponential envelope absorption.}
If the solution is known to have the form
$f^*(x) = g(x) \cdot e^{-\lambda x}$ where $\lambda > 0$ is known
(e.g., a dissipative system with known damping rate) and $g$ is an
unknown modulation, define:
\begin{equation}
  F(x, y, z) = y \cdot e^{\lambda x} - z.
\end{equation}
Then $h^*(x) = g(x)$ and $\beta = 1$.
The network learns only the slowly varying modulation $g$ instead
of the oscillatory, decaying product $f^* = g \cdot e^{-\lambda x}$.
Since the exponential envelope concentrates spectral energy at high
frequencies, removing it reduces the Barron norm:
$C_{h^*} = C_g \ll C_{g \cdot e^{-\lambda x}} = C_{f^*}$.

This explains why Laguerre-type networks (using $t^n e^{-\lambda t}$
as activation functions) outperform generic MLPs for damped systems:
the activation absorbs the exponential structure, and the weights
modulate a simpler residual.

\textbf{Case 2: Polynomial factor extraction.}
If $f^*(x) = p(x) \cdot s(x)$ where $p$ is a known polynomial
(e.g., from a Taylor expansion of the nonlinearity at a known
equilibrium) and $s$ is unknown but smooth, define:
\begin{equation}
  F(x, y, z) = y / p(x) - z.
\end{equation}
Then $h^*(x) = s(x)$ and $\beta = 1$.
This is a direct instance of Proposition~\ref{prop:factorizacion}.

\textbf{Case 3: Frequency demodulation.}
If the solution is known to oscillate at a dominant frequency
$\omega$ (e.g., an AC circuit, a mechanical resonator), the signal
can be decomposed as
$f^*(x) = A(x)\cos(\omega x) + B(x)\sin(\omega x)$
where $A, B$ are slowly varying envelopes.
Defining two implicit models:
\begin{align}
  F_c(x, y, z) &= \langle y, \cos(\omega \cdot)\rangle_{\mathrm{local}} - z, \\
  F_s(x, y, z) &= \langle y, \sin(\omega \cdot)\rangle_{\mathrm{local}} - z,
\end{align}
where $\langle\cdot,\cdot\rangle_{\mathrm{local}}$ denotes a
windowed inner product, the learning targets become $A$ and $B$
respectively, each with $\beta = 1$.

This is the mechanism underlying SIREN \cite{Sitzmann2020}:
the sinusoidal activation encodes the carrier frequency, and the
network weights learn the slowly varying modulation.
The framework explains why SIREN outperforms generic MLPs for
wave-like dynamics: the activation absorbs the dominant oscillatory
structure (reducing $B_{\mathrm{imp}}$) without amplification
($\beta = 1$).

\begin{remark}[Common structure of the three cases]
\label{rem:common_structure}
All three cases share the form $F(x,y,z) = \psi(x,y) - z$ with
$\beta = 1$.
The known structure (exponential envelope, polynomial factor, carrier
frequency) is absorbed by $\psi$, reducing the learning target to
a simpler modulation.
The structural content is $\mathcal{S}(F) = B_{\mathrm{exp}} /
B_{\mathrm{imp}} > 1$, and the gain condition~\eqref{eq:condicion_S}
reduces to $\mathcal{S}(F) > M_{\mathrm{imp}} / M_{\mathrm{exp}}$,
which is satisfied whenever the modulation has comparable range to
the original signal.
\end{remark}

\subsection{Homotopy Analysis Method}
\label{sec:ham}

The Homotopy Analysis Method (HAM) of Liao \cite{Liao2003, Liao2012}
provides a different and deeper instance of the framework.
Rather than absorbing a single known structural feature (an
exponential, a polynomial, a frequency), the HAM systematically
decomposes the \emph{entire} solution of a nonlinear ODE
$\mathcal{N}[u] = 0$ into a series of linear corrections.

The deformation is governed by the homotopy equation:
\begin{equation}
  (1-q)\,\mathcal{L}[u(t;q) - u_0(t)]
  = q\,\hbar\,H(t)\,\mathcal{N}[u(t;q)],
  \label{eq:HAM}
\end{equation}
where $\mathcal{L}$ is a chosen auxiliary linear operator,
$\hbar \ne 0$ is a convergence-control parameter,
$H(t)$ is an auxiliary function, and $q \in [0,1]$ is the homotopy
parameter.
Setting $q = 0$ gives $u(t;0) = u_0(t)$; setting $q = 1$ gives
$\mathcal{N}[u(t;1)] = 0$, the original equation.

The solution expands as:
\begin{equation}
  u(t;q) = u_0(t) + \sum_{k=1}^{\infty} u_k(t)\,q^k,
\end{equation}
where each $u_k$ satisfies a \emph{linear} equation with forcing
terms depending on $u_0, \ldots, u_{k-1}$.

In the language of the present framework:
the target function $f^*$ is the solution of $\mathcal{N}[u] = 0$
(potentially highly complex);
the function $h^*$ that the MLP must learn is the residual
correction beyond the $K$-th order HAM approximation:
$h^*(t) = u(t;1) - \sum_{k=0}^{K-1} u_k(t)$.
Since the $K$ leading terms are computed analytically by solving
linear equations, only the residual $h^*$ requires neural
approximation.

\begin{proposition}[HAM complexity reduction --- conditional]
\label{prop:HAM}
If the HAM series converges geometrically with ratio $\rho < 1$
(i.e., $\norm{u_k}_\infty \le C\rho^k$), and the Barron norm of
$h^*_K := u^* - \sum_{k=0}^{K-1} u_k$ satisfies
$C_{h^*_K} \le C' \rho^K$ for some constant $C'$, then the path norm
required to approximate $h^*_K$ satisfies:
\begin{equation}
  B_{\mathrm{imp}} = O(\rho^K) \ll B_{\mathrm{exp}}
\end{equation}
for $K$ sufficiently large, and the sample complexity ratio
decreases exponentially with the truncation order:
\begin{equation}
  \frac{N_{\mathrm{imp}}}{N_{\mathrm{exp}}}
  = O(\rho^{2K}).
\end{equation}
\end{proposition}

\begin{remark}[Scope of Proposition~\ref{prop:HAM}]
\label{rem:ham_scope}
The condition $C_{h^*_K} \le C'\rho^K$ is an additional regularity
assumption beyond $\norm{u_k}_\infty \le C\rho^k$.
The supremum norm and the Barron norm are not equivalent: a function
with small $\norm{\cdot}_\infty$ may have high spectral variation
(e.g., small-amplitude but highly oscillatory functions).
However, for the HAM residuals $h^*_K$, which are solutions of
\emph{linear} ODEs with exponentially decaying forcing, the spectral
content is controlled by the forcing spectrum, which itself decays.
A rigorous proof of $C_{h^*_K} = O(\rho^K)$ for specific operator
classes (e.g., $\mathcal{L}[u] = \ddot{u} + \omega^2 u$) requires
explicit Fourier analysis of the HAM recurrence and is left to
future work.
The proposition is stated conditionally on this spectral decay, which
is expected to hold for the standard classes of HAM-solvable ODEs.
\end{remark}

\begin{remark}[HAM vs.\ variable transformations]
\label{rem:ham_vs_vt}
The variable transformations of Section~\ref{sec:variable_transform}
require the practitioner to identify a single structural feature
(exponential, polynomial, frequency) and encode it manually in $F$.
The HAM provides a \emph{systematic} procedure: given only the
operator $\mathcal{N}$ and an auxiliary linear operator
$\mathcal{L}$, it generates a sequence of linear corrections that
progressively absorb the complexity of $f^*$.
No manual identification of structural features is required ---
the method constructs them automatically from the physics.
The price is the computation of $K$ terms of the homotopy series,
but the reward is an exponential reduction in sample complexity with
increasing $K$.
\end{remark}

%% =========================================================
\section{Interpretation: Structural Inductive Bias and Effective
  Hypothesis Space Reduction}
\label{sec:bias}
%% =========================================================

\subsection{Structural inductive bias}

In the explicit formulation $y = \MLP(x;\theta)$, the network must
discover \emph{all} structure of $f^*$ from data.
The only inductive bias is architectural (depth, width, activation).

In the implicit formulation $F(x, y, \MLP(x;\theta)) = 0$, the
function $F$ acts as an \emph{external structural inductive bias}:
it encodes prior knowledge about the relationship between $x$ and $y$
that does not need to be learned.
This knowledge may come from physical laws, geometric constraints,
symmetries, or analytical decompositions (such as the HAM series).

The distinction is analogous to the difference between a
noninformative and an informative prior in Bayesian inference: the
informative prior (when correct) concentrates the posterior more
rapidly, requiring fewer observations.
In our framework, the ``prior'' is the structure of $F$, and the
``concentration of the posterior'' manifests as reduction of the
Rademacher complexity.

\subsection{Effective hypothesis space reduction}

The nominal hypothesis class $\mathcal{H}_B$ is the same in both
formulations (same network, same weights).
What changes is the \emph{effective hypothesis space}: the subset of
$\mathcal{H}_B$ relevant to the problem.

In the explicit formulation, the network must cover all functions
$\varepsilon_0$-close to $f^*$:
\begin{equation}
  \mathcal{H}_{\mathrm{exp}}^{\mathrm{eff}}
  = \{g \in \mathcal{H}_{B_{\mathrm{exp}}} :
    \sup_x |g(x) - f^*(x)| \le \varepsilon_0\}.
\end{equation}

In the implicit formulation, the network only needs to cover
functions $\varepsilon_0$-close to $h^*$, which is simpler:
\begin{equation}
  \mathcal{H}_{\mathrm{imp}}^{\mathrm{eff}}
  = \{g \in \mathcal{H}_{B_{\mathrm{imp}}} :
    \sup_x |g(x) - h^*(x)| \le \varepsilon_0\}.
\end{equation}

Since $B_{\mathrm{imp}} < B_{\mathrm{exp}}$, the effective implicit
space has lower ``volume'' in the sense of Rademacher complexity.
The network retains its expressive capacity, but the implicit model
\emph{restricts the search} to a smaller region of parameter space.

\begin{remark}[Computational Occam principle]
The implicit formulation implements a version of Occam's principle:
among all networks satisfying
$F(x, y, \MLP(x;\theta)) = 0$, those with smaller path norm $B$
generalize better.
The function $F$ acts as a filter that eliminates hypotheses
incompatible with the known structure, reducing effective complexity
without reducing the network's expressive capacity.
\end{remark}

\begin{remark}[Connection to PAC-Bayes]
The effective space reduction argument is compatible with the
PAC-Bayes framework \cite{McAllester1999}, where the generalization
bound depends on the KL divergence between posterior and prior
distributions over $\theta$.
The function $F$ acts as a deterministic constraint that concentrates
the posterior mass on
$\Theta^*_{\mathrm{imp}} \subset \Theta$, reducing the effective KL
divergence.
Note that $F$ is not a prior in the strict sense (it is not a
probability distribution over $\theta$), but a deterministic
constraint that reduces the effective support.
The formalization of this connection within PAC-Bayes is left as
future work.
\end{remark}

\subsection{Why PINNs do not fit this framework}
\label{sec:why_not_pinns}

It is instructive to examine why the PINN approach
\cite{Raissi2019} does \emph{not} constitute an instance of the
present framework, despite sharing the goal of injecting physical
knowledge into neural models.

In a standard PINN, the network $u_\theta(x)$ approximates the PDE
solution $f^*$ \emph{directly}, and the loss function includes both
a data-fitting term and a PDE residual:
$\mathcal{L} = \mathcal{L}_{\mathrm{data}} +
\lambda\,\mathcal{L}_{\mathrm{PDE}}$.
The PDE term constrains the optimization landscape (it penalizes
networks incompatible with the governing equation), but the learning
target remains $f^*$ itself.
In the notation of this paper: $h^* = f^*$, so
$B_{\mathrm{imp}} = B_{\mathrm{exp}}$ and $\mathcal{S}(F) = 1$.

The PINN mechanism is \emph{regularization}: it reduces the
effective search space during optimization by penalizing physically
inconsistent solutions.
The mechanism formalized in Theorem~\ref{teo:principal} is
\emph{target transformation}: it changes what the network must
learn, replacing $f^*$ with a genuinely simpler function $h^*$.

This distinction has a precise consequence.
PINN-type regularization may improve optimization convergence and
reduce overfitting, but it does not change the Rademacher complexity
of the hypothesis class (the class of functions the network
\emph{can} represent is unchanged).
Target transformation reduces the Rademacher complexity because the
approximation target itself is simpler (lower Barron norm, lower path
norm).

The two mechanisms are complementary: one could use a target
transformation (absorbing known structure in $F$) together with a
PINN-style residual loss on the transformed problem, obtaining the
benefits of both.

%% =========================================================
\section{Discussion and Future Directions}
\label{sec:discussion}
%% =========================================================

\subsection{What this paper proves}

Theorem~\ref{teo:principal} establishes that, conditional on
Assumption~\ref{hip:reduccion}, the implicit model
$F(x, y, \MLP(x;\theta)) = 0$ achieves a tighter generalization
bound than the explicit model $y = \MLP(x;\theta)$.
Corollary~\ref{cor:datos} translates this into an explicit sample
complexity ratio.
The gain condition~\eqref{eq:condicion_S} is a verifiable criterion
that balances the path norm reduction against the Lipschitz
amplification introduced by $F$.

The result is constructive: the hypotheses are stated explicitly,
the conditions are verifiable for specific choices of $F$, and the
counterexamples show that the result is tight (not every implicit
model helps).

\subsection{What this paper does not prove}

The paper does not resolve two important questions:

\textbf{(i) When does Assumption~\ref{hip:reduccion} hold?}
The sufficient conditions of Section~\ref{sec:sufficient} cover
two important classes (structural factorization and nonlinearity
absorption), but do not exhaust all possibilities.
A general characterization of which functions $F$ yield
$B_{\mathrm{imp}} < B_{\mathrm{exp}}$ for a given class of target
functions would require a deeper analysis of the interaction between
the spectral structure of $f^*$ and the transformation induced
by $F$.

\textbf{(ii) Why do certain activation functions outperform others
for specific dynamical classes?}
The framework shows that structural knowledge in $F$ reduces sample
complexity, but does not explain at a deeper level why
$\sin(\omega t)$ is the ``right'' activation for oscillatory systems
or $e^{-\lambda t}$ for damped systems --- beyond the observation
that they lead to lower $B_{\mathrm{imp}}$.

\subsection{Toward autosimilar function spaces}

Both questions point toward a topological framework that we outline
here and defer to a companion paper.

The key observation is that the history of function approximation
can be read as a sequence of discoveries of \emph{autosimilar bases},
each adapted to a different class of natural structure:
Fourier series for oscillatory dynamics,
Laguerre functions for damped systems,
Chebyshev polynomials for compact-domain approximation,
Gaussian RBFs for spatially localized problems,
and wavelet bases for multi-scale phenomena.

Each of these bases is the orbit of a single \emph{primitive
function} under a family of contractive transformations ---
precisely the structure of an iterated function system (IFS) in the
sense of Hutchinson \cite{Hutchinson1981}.
The Homotopy Analysis Method \cite{Liao2003, Liao2012} provides the
missing link: it constructs the solution of a nonlinear ODE as the
limit of an orbit $\{v_k = \Gamma_\hbar^k[u_0]\}$ generated by a
contractive operator, where the primitive $u_0$ is the homogeneous
solution of the auxiliary linear operator $\mathcal{L}$.

Three consequences follow:
\begin{enumerate}
  \item The exact solution $u^*$ is the attractor of the IFS, not an
    external target.
    Approximation is an intrinsic property of the space the solution
    inhabits, not an operation imposed from outside.
  \item Every HAM scheme defines a neural architecture whose
    activation function is $u_0$ and whose basis is generated by the
    HAM contraction.
    Conversely, every neural architecture whose activation is the
    homogeneous solution of some $\mathcal{L}$ is an instance of a
    HAM scheme.
  \item The canonical neural architectures (SIREN, RBF, Chebyshev,
    wavelet networks) are the neural counterparts of the canonical
    HAM operators.
\end{enumerate}

These claims require precise definitions of autosimilar function
spaces, contractivity conditions on $\Gamma_\hbar$, and a formal
bridge theorem connecting HAM bases to neural architectures.
Their development and proof constitute the subject of a companion
paper \cite{Rodrigo2026companion}, in which the topological
framework provides the structural explanation for \emph{why}
Assumption~\ref{hip:reduccion} holds for physics-informed
architectures, complementing the statistical guarantee established
here.

%% =========================================================
\section{Summary of Results}
\label{sec:summary}
%% =========================================================

\begin{center}
\renewcommand{\arraystretch}{1.5}
\begin{tabular}{lc}
  \toprule
  \textbf{Result} & \textbf{Equation} \\
  \midrule
  Explicit generalization bound &
  $O\!\left(\dfrac{M_{\mathrm{exp}}\, B_{\mathrm{exp}}\, X_{\max}
    \sqrt{L\log d}}{\sqrt{N}}\right)$ \\[8pt]
  Implicit generalization bound &
  $O\!\left(\dfrac{\beta\, M_{\mathrm{imp}}\, B_{\mathrm{imp}}\,
    X_{\max} \sqrt{L\log d}}{\sqrt{N}}\right)$ \\[8pt]
  Sample complexity ratio &
  $\dfrac{\beta^2 M_{\mathrm{imp}}^2 B_{\mathrm{imp}}^2}
    {M_{\mathrm{exp}}^2 B_{\mathrm{exp}}^2}$ \\[8pt]
  Structural content $\mathcal{S}(F)$ &
  $B_{\mathrm{exp}} / B_{\mathrm{imp}}$ \\[4pt]
  Gain condition &
  $\mathcal{S}(F) > \beta \cdot M_{\mathrm{imp}} / M_{\mathrm{exp}}$
  \\
  \bottomrule
\end{tabular}
\end{center}

%% =========================================================
\section{Conclusions}
\label{sec:conclusions}
%% =========================================================

We have established that implicit neural models with structural
inductive bias achieve tighter generalization bounds than explicit
models, under quantitative conditions that balance the path norm
reduction against the Lipschitz amplification introduced by the
implicit function $F$.

The essential mechanism is \emph{target transformation}: the
implicit function $F$ changes the learning target from the original
$f^*$ to a simpler residual $h^*$, reducing the intrinsic
approximation complexity.
This mechanism is fundamentally different from the regularization
approach of PINNs, which constrains the optimization landscape but
does not change the learning target.

Two instantiations have been identified:
\begin{enumerate}
  \item \textbf{Physics-informed variable transformations}, where
    known structural features (exponential envelopes, polynomial
    factors, carrier frequencies) are absorbed by $F$ with
    $\beta = 1$ and maximal structural content.
  \item \textbf{Homotopy Analysis Method decompositions}, where the
    HAM series provides a systematic decomposition of the target
    function into linear corrections of geometrically decreasing
    complexity, requiring no manual identification of structural
    features.
\end{enumerate}

The practical consequence is a principled criterion for evaluating
whether a given implicit formulation will outperform the explicit
baseline: compute the structural content $\mathcal{S}(F)$ and verify
the gain condition $\mathcal{S}(F) > \beta \cdot M_{\mathrm{imp}} /
M_{\mathrm{exp}}$.

The deeper question --- \emph{why} certain bases are intrinsically
adapted to certain dynamics --- requires a topological framework
based on autosimilar function spaces and the bridge between HAM and
neural architectures, to be developed in a companion paper.

\section*{Declaration of competing interest}
The author declares no competing interests.

\section*{Acknowledgments}
This work was supported by CONICET and Universidad Nacional de San
Juan, Argentina.

\bibliographystyle{plain}
\begin{thebibliography}{99}

\bibitem{Cybenko1989}
G.~Cybenko,
Approximation by superpositions of a sigmoidal function,
\emph{Math.\ Control Signals Syst.} 2~(4) (1989) 303--314.

\bibitem{Hornik1991}
K.~Hornik,
Approximation capabilities of multilayer feedforward networks,
\emph{Neural Netw.} 4~(2) (1991) 251--257.

\bibitem{Funahashi1989}
K.-I.~Funahashi,
On the approximate realization of continuous mappings by neural
networks,
\emph{Neural Netw.} 2~(3) (1989) 183--192.

\bibitem{Barron1993}
A.~R.~Barron,
Universal approximation bounds for superpositions of a sigmoidal
function,
\emph{IEEE Trans.\ Inform.\ Theory} 39~(3) (1993) 930--945.

\bibitem{Bartlett2002}
P.~L.~Bartlett, S.~Mendelson,
Rademacher and Gaussian complexities: risk bounds and structural
results,
\emph{J.~Mach.\ Learn.\ Res.} 3 (2002) 463--482.

\bibitem{Golowich2018}
N.~Golowich, A.~Rakhlin, O.~Shamir,
Size-independent sample complexity of neural networks,
in: \emph{Proc.\ COLT}, 2018, pp.~297--299.

\bibitem{Ledoux1991}
M.~Ledoux, M.~Talagrand,
\emph{Probability in Banach Spaces}, Springer, 1991.

\bibitem{Neyshabur2015}
B.~Neyshabur, R.~Tomioka, N.~Srebro,
Norm-based capacity control in neural networks,
in: \emph{Proc.\ COLT}, 2015.

\bibitem{Raissi2019}
M.~Raissi, P.~Perdikaris, G.~E.~Karniadakis,
Physics-informed neural networks: a deep learning framework for
solving forward and inverse problems involving nonlinear partial
differential equations,
\emph{J.~Comput.\ Phys.} 378 (2019) 686--707.

\bibitem{Chen2018}
R.~T.~Q.~Chen, Y.~Rubanova, J.~Bettencourt, D.~Duvenaud,
Neural ordinary differential equations,
in: \emph{Adv.\ Neural Inform.\ Process.\ Syst.}, vol.~31, 2018.

\bibitem{Bai2019}
S.~Bai, J.~Z.~Kolter, V.~Koltun,
Deep equilibrium models,
in: \emph{Adv.\ Neural Inform.\ Process.\ Syst.}, vol.~32, 2019.

\bibitem{Sitzmann2020}
V.~Sitzmann, J.~N.~P.~Martel, A.~W.~Bergman, D.~B.~Lindell,
G.~Wetzstein,
Implicit neural representations with periodic activation functions,
in: \emph{Adv.\ Neural Inform.\ Process.\ Syst.}, vol.~33, 2020,
pp.~7462--7473.

\bibitem{Liao2003}
S.-J.~Liao,
\emph{Beyond Perturbation: Introduction to the Homotopy Analysis
Method}, Chapman \& Hall/CRC, 2003.

\bibitem{Liao2012}
S.-J.~Liao,
\emph{Homotopy Analysis Method in Nonlinear Differential Equations},
Springer, 2012.

\bibitem{Hutchinson1981}
J.~E.~Hutchinson,
Fractals and self-similarity,
\emph{Indiana Univ.\ Math.\ J.} 30~(5) (1981) 713--747.

\bibitem{Vapnik1995}
V.~N.~Vapnik,
\emph{The Nature of Statistical Learning Theory}, Springer, 1995.

\bibitem{Mitchell1980}
T.~M.~Mitchell,
The need for biases in learning generalizations,
Technical Report CBM-TR-117, Rutgers University, 1980.

\bibitem{McAllester1999}
D.~A.~McAllester,
PAC-Bayesian model averaging,
in: \emph{Proc.\ COLT}, 1999, pp.~164--170.

\bibitem{Rodrigo2026companion}
R.~O.~Rodrigo,
Physics-driven autosimilar function spaces via the Homotopy Analysis
Method: a unifying framework for neural approximation,
2026, in preparation.

\end{thebibliography}

\end{document}
